\documentclass[11pt,a4paper]{article}

\usepackage[margin=1in]{geometry}
\usepackage{amsmath,amssymb,amsthm}
\usepackage{graphicx}
\usepackage{booktabs}
\usepackage{microtype}
\usepackage{xcolor}
\usepackage[hidelinks,colorlinks=true,linkcolor=blue!50!black,citecolor=blue!40!black,urlcolor=blue!50!black]{hyperref}
\usepackage[numbers,sort&compress]{natbib}

\newtheorem{theorem}{Theorem}
\newtheorem{proposition}[theorem]{Proposition}
\newtheorem{lemma}[theorem]{Lemma}
\newtheorem{corollary}[theorem]{Corollary}
\theoremstyle{definition}
\newtheorem{definition}[theorem]{Definition}
\newtheorem{conditional}[theorem]{Conditional Proposition}
\theoremstyle{remark}
\newtheorem{remark}[theorem]{Remark}

\title{\textbf{Classical Primes as Closure-Irreducibles}\\[2pt]
\large The $L_4$ $\sigma$-Classification via $\mathbb{Q}(\sqrt{5})$ and the Universal $\sigma$-Pattern Across the Closure-Irreducibility Typology\\[2pt]
\normalsize \emph{Closure-Irreducibles, Paper II of a paired contribution.}\\[1pt]
\small Companion: \emph{The Critical Line as Observation Manifold: A Closure-Theoretic Foundation} (Paper~I).\\[1pt]
\small (Open research note --- pre-peer-review draft)}

\author{%
  Lee Smart\\[4pt]
  \textit{Institute of Vibrational Field Dynamics}\\[2pt]
  \texttt{contact@vibrationalfielddynamics.org}\\[2pt]
  \texttt{@vfd\_org}%
}
\date{May 2026}

\begin{document}
\maketitle
\thispagestyle{empty}

\begin{abstract}
\noindent
\textbf{Scope.} This paper develops the $L_4$ level of the closure-irreducibility
typology introduced in Paper~I~\cite{CriticalLinePaperI} --- namely, the
placement of classical integer primes within the cascade $\sigma$-structure.
\textbf{Non-claim.} It does \emph{not} claim to derive classical primes
from cascade closure. The integers $2, 3, 5, 7, 11, \ldots$ exist
independently as facts of arithmetic; we are not redefining them.
\textbf{Method.} We use the classical Dedekind splitting law in the
cascade's natural ring of integers $\mathbb{Z}[\varphi]$
(unconditional, well-known since 1871) to demonstrate that every
classical prime falls into exactly one of three cascade-$\sigma$-classes
under the Galois conjugation $\sqrt{5} \mapsto -\sqrt{5}$.
\textbf{Primary structural result.} The $L_4$ $\sigma$-classification
of classical integer primes (split / inert / ramified) is the natural
manifestation, at the integer level, of the same $\sigma$-pattern that
governs $L_0$ through $L_3$ in the closure-irreducibility typology
(vertices, $\sigma$-orbits, spectral eigenspaces, and pentagonal-clock
$K$-poles). \textbf{Universal $\sigma$-pattern theorem.} Every level
$L_0$--$L_4$ of the typology admits a $\sigma$-classification of the
same shape: a fixed subset, a paired subset, and (in some levels) a
single ramified case. The $\sigma$-pattern is universal across all
five levels.
\textbf{What this provides.} An explicit structural placement of
classical primes within the cascade typology --- demonstrating, under
the named hypothesis H-SigmaRestrict (the identification of the cascade
spectral involution with the Galois conjugation on $\mathbb{Q}(\sqrt{5})$),
that they constitute one specific level of the closure-irreducibility
typology. The placement is unconditional in the underlying mathematics
(classical Dedekind splitting law) and structural in its cascade
interpretation.
\textbf{What this does not provide.} A derivation of classical primes
from cascade machinery; a bridge from cascade's analytic object
$\hat{\zeta}(z)$ to classical $\zeta(s)$ at the level of generating
functions (blocked by the Hecke obstruction, \emph{rh-formal}
$M3$~\cite{rhFormalSmart}); a proof, conditional or otherwise, of the
classical Riemann Hypothesis.
\textbf{Relation to classical number theory.} The Dedekind splitting law
in $\mathbb{Z}[\varphi]$ is classical, unconditional, and has been
textbook material for over a century. The contribution of the present
paper is its \emph{reinterpretation} as the $L_4$ manifestation of the
cascade $\sigma$-typology, not its discovery.
\end{abstract}

\bigskip

\noindent\fbox{%
\parbox{0.96\linewidth}{%
\textbf{Plain-language hypothesis.}

\smallskip\noindent
Classical primes $(2, 3, 5, 7, 11, \ldots)$ remain classical primes.
We are not redefining them. We are showing that they fit naturally into
the closure-irreducibility typology introduced in Paper~I, as one
specific level ($L_4$).

\smallskip\noindent
The placement is unconditional and classical: every prime $p$ has a
precise behaviour under the Galois conjugation $\sqrt{5} \mapsto -\sqrt{5}$
in the cascade's natural number field $\mathbb{Q}(\sqrt{5})$. Determined
by $p \bmod 5$:
\begin{itemize}\itemsep0pt
  \item $p \equiv 1$ or $4 \pmod{5}$: \textbf{splits} (e.g., $11, 19, 29, 31, \ldots$) --- $\sigma$-paired pair
  \item $p \equiv 2$ or $3 \pmod{5}$: \textbf{inert} (e.g., $2, 3, 7, 13, \ldots$) --- $\sigma$-fixed
  \item $p = 5$: \textbf{ramified} --- single fixed point
\end{itemize}

\smallskip\noindent
This $\sigma$-classification has \emph{the same shape} as the
$\sigma$-classification at every other level of the typology
($L_0$ vertices, $L_1$ orbits, $L_2$ spectral classes, $L_3$ $K$-poles).
The $\sigma$-pattern --- fixed subset, paired subset, possibly a single
ramified case --- is universal across all five levels.

\smallskip\noindent
What is new is not the splitting law (Dedekind, 1871). What is new is
\emph{the recognition that this classical law is the $L_4$ instance of
a universal cascade $\sigma$-pattern}.
}}

\bigskip

\paragraph{Programme map.}
This paper sits in the closure-foundations layer of the Existence /
Life / Closure programme. It is Paper~II of the Closure-Irreducibles
paired contribution. Paper~I~\cite{CriticalLinePaperI} introduces the
five-level typology and proves the critical-line / observation-manifold
foundation. The present paper develops the $L_4$ level in detail and
demonstrates the universal $\sigma$-pattern across all five levels.
Reading either paper alone is self-contained; reading both gives the
complete argument.

\paragraph{Status taxonomy.}
\textbf{Theorem} = formal structural claim under stated assumptions.
\textbf{Classical Theorem} = result imported from classical number theory,
cited with reference (e.g., Dedekind splitting law).
\textbf{Conditional Proposition} = bridge claim depending on a named
hypothesis. \textbf{Foundational Interpretation} = structural reading
of classical mathematical content within the closure framework, marked
as such.

\paragraph{Non-claims (load-bearing).}
This paper does \textbf{not} claim that:
\begin{itemize}\itemsep0pt
  \item Classical integer primes are derived from cascade closure (they are not; they exist as facts of arithmetic).
  \item The classical Riemann Hypothesis is proved or addressed at the arithmetic level (it is not; full derivation of $\zeta$ from cascade is blocked by the Hecke obstruction \cite{rhFormalSmart}).
  \item The definition of ``integer prime'' is changed (it is not; classical primality is exactly what it has always been).
  \item Classical number theory is wrong about anything (it is not; this paper reinterprets, it does not correct).
  \item Cascade's analytic object $\hat{\zeta}(z)$ equals or replaces classical $\zeta(s)$ (it does not; they are distinct analytic objects with different multiplicity structures).
\end{itemize}
The framework is one of \emph{reinterpretation through the closure
lens}, not of redefinition.

\section{Introduction: the bridge question}\label{sec:intro}

Paper~I established that the critical line $\mathrm{Re}(s) = 1/2$ is the
$\sigma$-fixed set of an involution that closure naturally produces on
$V_{600}$, and introduced a five-level typology of
\emph{closure-irreducibles} ($L_0$--$L_4$). Within that typology,
classical integer primes ($L_4$) appeared as one declared level but were
not developed in detail.

The natural next question is: \emph{how, exactly, do classical primes
fit into the cascade $\sigma$-typology?}

There are three logically distinct ways to answer:

\begin{enumerate}\itemsep0pt
  \item \textbf{Derivation.} Show that the integers $2, 3, 5, 7, 11, \ldots$
    arise from cascade machinery (e.g., as the spectrum of some
    cascade-natural operator, or as the support of some
    cascade-natural distribution).
  \item \textbf{Classification.} Show that the classical primes, taken
    as given, fit naturally into the cascade $\sigma$-typology via some
    structural action of cascade on $\mathbb{Z}$.
  \item \textbf{Identification.} Show that cascade's analytic object
    $\hat{\zeta}(z)$ equals or relates to classical $\zeta(s)$ in a way
    that makes their content identifiable.
\end{enumerate}

The strongest answer would be all three. Currently:

\begin{itemize}\itemsep0pt
  \item Option~3 (\textbf{identification}) is \emph{blocked} by the
    Hecke obstruction (rh-formal $M3$~\cite{rhFormalSmart}): cascade's
    $\hat{\zeta}(z)$ has multiplicities $(1, 1, 5, 5)$ at four
    pentagonal-clock $K$-values, whereas classical $\zeta(s)$ has
    multiplicities all $1$ at infinitely many primes. These analytic
    objects do not align.
  \item Option~1 (\textbf{derivation}) is \emph{open} and likely deep:
    classical primes have an infinite, irregular distribution
    governed by analytic number theory; cascade closure operates on a
    finite polytope $V_{600}$, and bridging finite cascade content to
    infinite arithmetic is non-trivial.
  \item Option~2 (\textbf{classification}) is \emph{available now},
    unconditionally, using classical number theory that has been known
    since Dedekind's work in 1871.
\end{itemize}

This paper completes Option~2.

\section{The cascade's natural number field: $\mathbb{Q}(\sqrt{5})$ and $\mathbb{Z}[\varphi]$}\label{sec:Q_sqrt5}

We begin by recording what the cascade's natural arithmetic substrate is.

\subsection{Why $\mathbb{Q}(\sqrt{5})$?}

The closure framework on $V_{600}$ is built over the golden-ratio field.
Concretely:

\begin{itemize}\itemsep0pt
  \item The $600$-cell has edge length $1 / \varphi$, where $\varphi = (1 + \sqrt{5})/2$.
  \item The adjacency operator $A_{V_{600}}$ has eigenvalue $+6\varphi$ at the $\sigma$-paired dipole class (Paper~I, Theorem~$T1$).
  \item The closure operator is $\mathcal{C}_\varphi = L_{V_{600}} + \varphi^{-2} I$.
  \item The $120$ vertices of $V_{600}$ form the binary icosahedral group $2I$, and lie naturally in the icosian quaternion ring built over $\mathbb{Z}[\varphi]$.
\end{itemize}

Hence:

\begin{definition}[Cascade's natural number field]\label{def:natural_field}
The cascade's natural number field is $K = \mathbb{Q}(\sqrt{5})$, with
ring of integers $\mathbb{Z}[\varphi] = \mathbb{Z}[(1+\sqrt{5})/2]$.
\end{definition}

\subsection{The cascade $\sigma$-action restricted to $\mathbb{Q}(\sqrt{5})$}

\begin{definition}[Galois conjugation on $K$]\label{def:galois}
The non-trivial element of $\mathrm{Gal}(K/\mathbb{Q})$ is the involution
$\sigma_K : \sqrt{5} \mapsto -\sqrt{5}$, equivalently
$\varphi \mapsto 1 - \varphi = -1/\varphi$.
\end{definition}

\begin{conditional}[Cascade $\sigma$ restricts to $\sigma_K$, under H-SigmaRestrict]\label{lem:sigma_restricts}
\emph{Assume H-SigmaRestrict} --- the structural identification, made
in \cite{SmartExistenceClosure} (Paper~I of the existence programme)
and \cite{rhFormalSmart} (Section~$2$), that the cascade spectral
involution $\tau_{\mathrm{spec}}$ on $V_{600}$ acts on icosian quaternion
coordinates via $\varphi \mapsto -1/\varphi$ on the underlying
coefficient field. Then $\tau_{\mathrm{spec}}$, restricted to coefficients
in $\mathbb{Q}(\sqrt{5})$, induces the Galois conjugation
$\sigma_K : \sqrt{5} \mapsto -\sqrt{5}$.
\end{conditional}

\begin{proof}[Derivation under H-SigmaRestrict]
By definition of $\tau_{\mathrm{spec}}$ (Paper~I), the spectral
involution negates the $+6\varphi$ eigenspace of $A_{V_{600}}$, which
under H-SigmaRestrict corresponds to $\varphi \mapsto -1/\varphi$ on
the icosian coefficient field. Since
$-1/\varphi = 1 - \varphi = (1 - \sqrt{5})/2$ and
$\varphi = (1 + \sqrt{5})/2$, the action $\varphi \mapsto -1/\varphi$
sends $\sqrt{5}$ to $-\sqrt{5}$, which is precisely
$\sigma_K \in \mathrm{Gal}(\mathbb{Q}(\sqrt{5}) / \mathbb{Q})$.
\end{proof}

\begin{remark}[Status of H-SigmaRestrict]\label{rmk:status_sigma}
H-SigmaRestrict is the icosian-quaternion construction of
\cite{SmartExistenceClosure} (Paper~I of the existence programme),
worked out in rh-formal \cite{rhFormalSmart} Section~$2$. It is
\emph{not} re-derived in the present paper. A reader who treats
H-SigmaRestrict as an open structural identification should regard the
$L_4$ placement (Conditional~\ref{thm:L4_placement} below) as
conditional accordingly. Independently of H-SigmaRestrict, the
classical Dedekind splitting law (Theorem~\ref{thm:dedekind}) holds
unconditionally; what H-SigmaRestrict provides is the
\emph{identification} of the cascade-side involution with the classical
Galois action. Without this identification the splitting law remains
classical mathematics; with it, the splitting becomes the $L_4$ instance
of the cascade $\sigma$-pattern.
\end{remark}

Under H-SigmaRestrict, we have the structural link: \emph{the cascade
$\sigma$ on $V_{600}$ and the classical $\sigma_K$ on
$\mathbb{Q}(\sqrt{5})$ are the same involution viewed in different
coordinates.} The remainder of the paper assumes H-SigmaRestrict to
state the $L_4$ placement and the universal pattern; the alternative
unconditional reading (without H-SigmaRestrict) is recorded in the
falsifier-path table at the end.

\section{The classical Dedekind splitting law in $\mathbb{Z}[\varphi]$}\label{sec:dedekind}

We import the classical splitting law, unmodified.

\begin{theorem}[Dedekind splitting law in the ring of integers of $\mathbb{Q}(\sqrt{5})$]\label{thm:dedekind}
Let $p$ be a rational prime. Then in $\mathbb{Z}[\varphi]$:
\begin{itemize}\itemsep0pt
  \item If $p \equiv 1$ or $4 \pmod{5}$, then $p$ \textbf{splits}:
    $p = \pi \bar{\pi}$ for some prime $\pi \in \mathbb{Z}[\varphi]$,
    with $\bar{\pi} = \sigma_K(\pi)$ the Galois conjugate.
  \item If $p \equiv 2$ or $3 \pmod{5}$, then $p$ is \textbf{inert}:
    $p$ remains prime in $\mathbb{Z}[\varphi]$.
  \item If $p = 5$, then $p$ is \textbf{ramified}:
    $5 = (\sqrt{5})^2 \cdot u$ for a unit $u \in \mathbb{Z}[\varphi]^\times$.
\end{itemize}
\end{theorem}

\begin{proof}[Proof (classical)]
This is the standard Dedekind splitting law in a quadratic number field,
applied to $K = \mathbb{Q}(\sqrt{5})$ with discriminant $5$. The
behaviour of $p$ in $\mathbb{Z}[\varphi]$ is determined by the Kronecker
symbol $(p/5)$ via quadratic reciprocity: split when $(p/5) = +1$
($p \equiv \pm 1 \pmod{5}$), inert when $(p/5) = -1$ ($p \equiv \pm 2 \pmod{5}$),
ramified when $p = 5$. See any standard algebraic-number-theory textbook
(Neukirch, Marcus, Stewart--Tall) for full proof.
\end{proof}

This is unconditional, classical mathematics dating to Dedekind 1871. We
do not reprove it. We import it.

\subsection{Examples}\label{ssec:examples}

Concrete classification of the small primes:

\begin{center}\small
\begin{tabular}{cccc}\toprule
$p$ & $p \bmod 5$ & Behaviour in $\mathbb{Z}[\varphi]$ & $\sigma$-class \\\midrule
$2$ & $2$ & inert & $\sigma$-fixed \\
$3$ & $3$ & inert & $\sigma$-fixed \\
$5$ & $0$ & \textbf{ramified}: $5 = (\sqrt{5})^2 \cdot u$ & single ramification \\
$7$ & $2$ & inert & $\sigma$-fixed \\
$11$ & $1$ & \textbf{splits}: $11 = \pi \bar{\pi}$, with $\pi = 3 + \varphi$ (one choice) & $\sigma$-paired pair \\
$13$ & $3$ & inert & $\sigma$-fixed \\
$17$ & $2$ & inert & $\sigma$-fixed \\
$19$ & $4$ & \textbf{splits} & $\sigma$-paired pair \\
$23$ & $3$ & inert & $\sigma$-fixed \\
$29$ & $4$ & \textbf{splits} & $\sigma$-paired pair \\
$31$ & $1$ & \textbf{splits} & $\sigma$-paired pair \\
$37$ & $2$ & inert & $\sigma$-fixed \\
$41$ & $1$ & \textbf{splits} & $\sigma$-paired pair \\
$43$ & $3$ & inert & $\sigma$-fixed \\
$47$ & $2$ & inert & $\sigma$-fixed \\
$53$ & $3$ & inert & $\sigma$-fixed \\
$59$ & $4$ & \textbf{splits} & $\sigma$-paired pair \\
$61$ & $1$ & \textbf{splits} & $\sigma$-paired pair \\
\bottomrule\end{tabular}
\end{center}

Every classical prime falls in exactly one of the three columns. There
are no exceptions, no edge cases.

\subsection{Asymptotic density}\label{ssec:density}

By Dirichlet's theorem on arithmetic progressions and the Chebotarev
density theorem, the asymptotic density of primes in each $\sigma$-class
is:

\begin{itemize}\itemsep0pt
  \item Split primes ($p \equiv \pm 1 \pmod{5}$): density $1/2$
  \item Inert primes ($p \equiv \pm 2 \pmod{5}$): density $1/2$
  \item Ramified prime ($p = 5$): density $0$ (finite set with one element)
\end{itemize}

Half the primes go to each of the two main $\sigma$-classes, with $5$
as the single ramification point. This is classical, unconditional.

\section{The $L_4$ $\sigma$-classification of classical primes}\label{sec:L4}

We now record the central definition and placement.

\begin{definition}[$L_4$ $\sigma$-classification]\label{def:L4_classification}
The \emph{$L_4$ $\sigma$-classification} of classical integer primes is
the partition of the prime set $\{p : p \text{ rational prime}\}$ into
three subsets determined by behaviour in $\mathbb{Z}[\varphi]$:
\begin{itemize}\itemsep0pt
  \item $\mathcal{P}_{\mathrm{split}} = \{p : p \text{ splits in } \mathbb{Z}[\varphi]\}$
  \item $\mathcal{P}_{\mathrm{inert}} = \{p : p \text{ inert in } \mathbb{Z}[\varphi]\}$
  \item $\mathcal{P}_{\mathrm{ram}} = \{5\}$
\end{itemize}
\end{definition}

\begin{conditional}[Placement of $L_4$ in cascade $\sigma$-typology, under H-SigmaRestrict]\label{thm:L4_placement}
Under Conditional~\ref{lem:sigma_restricts} (H-SigmaRestrict), the
$L_4$ $\sigma$-classification of Definition~\ref{def:L4_classification}
is the manifestation, at the level of classical integer primes, of the
cascade $\sigma$-action on $\mathbb{Q}(\sqrt{5})$. Specifically:
\begin{itemize}\itemsep0pt
  \item Split primes are $\sigma$-paired pairs $(\pi, \bar{\pi})$ with $\bar{\pi} = \sigma_K(\pi)$.
  \item Inert primes are $\sigma$-fixed (the principal ideal $(p)$ is preserved by $\sigma_K$).
  \item The ramified prime $p = 5$ is the unique fixed point of $\sigma_K$ acting on the prime ideals dividing the discriminant of $K$.
\end{itemize}

\noindent\emph{Unconditional content (independent of H-SigmaRestrict):} the
three-way classification of classical primes into split / inert /
ramified is the classical Dedekind splitting law
(Theorem~\ref{thm:dedekind}). H-SigmaRestrict supplies the
\emph{identification} that this three-way pattern is the cascade
$\sigma$-pattern; without H-SigmaRestrict the pattern remains classical
but does not formally constitute the $L_4$ instance of cascade
$\sigma$-action.
\end{conditional}

\begin{proof}[Derivation under H-SigmaRestrict]
Direct from Theorem~\ref{thm:dedekind} (classical splitting law,
unconditional) together with Conditional~\ref{lem:sigma_restricts}
(H-SigmaRestrict, identification of cascade $\sigma$ with $\sigma_K$).
Split primes factor as $\pi \bar{\pi}$ with $\bar{\pi} = \sigma_K(\pi)$
by definition; inert primes have $\sigma_K((p)) = (p)$ as the only
ideal of norm $p^2$; the ramified prime $5$ is the unique discriminant
divisor and is fixed by $\sigma_K$ up to units.
\end{proof}

Under H-SigmaRestrict, this places $L_4$ in the cascade $\sigma$-typology
with the same shape as $L_0$--$L_3$ in Paper~I: a fixed subset, a paired
subset, and (in this case) a single ramified prime.

\section{The universal $\sigma$-pattern across the typology}\label{sec:universal}

We can now state the headline structural observation of this paper.

\begin{conditional}[Universal $\sigma$-pattern, under H-SigmaRestrict $+$ H-UniqueSigma]\label{thm:universal}
Under Conditional~\ref{lem:sigma_restricts} (H-SigmaRestrict) and the
additional hypothesis H-UniqueSigma --- that at each level $L_i$ the
cascade $\sigma$ is the unique natural involution associated with the
closure operator $\mathcal{C}_\varphi$ on that level's objects, not
chosen by post-hoc selection --- every level $L_0, L_1, L_2, L_3, L_4$
of the closure-irreducibility typology admits a $\sigma$-classification
of the same structural shape: a $\sigma$-fixed subset $\mathcal{F}_L$,
a $\sigma$-paired subset $\mathcal{P}_L$, and (in some levels) a single
ramified or singular element $\mathcal{R}_L$. Concretely:
\end{conditional}

\begin{center}\small
\begin{tabular}{lp{0.18\linewidth}p{0.18\linewidth}p{0.18\linewidth}p{0.12\linewidth}}\toprule
Level & Object set & $\sigma$-fixed & $\sigma$-paired & Special \\\midrule
$L_0$ & $120$ $V_{600}$ vertices & vertices in $\sigma$-fixed orbits & vertices in $\sigma$-paired orbits & --- \\
$L_1$ & $107$ $\sigma$-orbits & $94$ fixed orbits & $13$ paired orbit pairs & --- \\
$L_2$ & $9$ spectral irreps & $\tau_{\mathrm{spec}}$-fixed eigenspaces & $\tau_{\mathrm{spec}}$-anti-fixed eigenspaces & --- \\
$L_3$ & $4$ $K$-poles & $K = 0$ (on critical line) & $K \in \{20, 52, 72\}$ (paired about line) & --- \\
$L_4$ & $\infty$ classical primes & $\mathcal{P}_{\mathrm{inert}}$ (e.g., $2, 3, 7, 13, \ldots$) & $\mathcal{P}_{\mathrm{split}}$ (e.g., $11, 19, 29, 31, \ldots$) & $\mathcal{P}_{\mathrm{ram}} = \{5\}$ \\
\bottomrule\end{tabular}
\end{center}

\begin{proof}[Proof sketch]
$L_0$, $L_1$: from Paper~I (substrate $\sigma$-orbit count $94 + 13 \cdot 2 = 120$).
$L_2$: from cascade spectral character theory (Paper~I, $W(H_4)$
character tables on $V_{600}$).
$L_3$: from rh-formal $M1$--$M2$ (pentagonal-clock closed form
$\hat{\zeta}(z) = \prod_K (1 - L_K z^{10} + z^{20})^{-m_K}$ with
$K \in \{0, 20, 52, 72\}$).
$L_4$: from Theorem~\ref{thm:L4_placement} above (Dedekind splitting
law applied through cascade $\sigma$).
\end{proof}

\subsection{What the universal pattern says}

Across all five levels, the same structural shape appears:

\begin{itemize}\itemsep0pt
  \item A set of closure-irreducibles at that level.
  \item Acted on by a $\sigma$-involution (Galois $\sigma_K$ on $\mathbb{Q}(\sqrt{5})$ at the integer level, cascade $\sigma$ at the geometric / spectral levels).
  \item Producing a partition into fixed elements, paired elements, and (at some levels) a single ramified element.
\end{itemize}

The $\sigma$-pattern is \emph{not} the same numerical count at each
level (the counts vary: $\{94, 13\}$ at $L_1$, $\{1, 1, 5, 5\}$ at $L_3$,
infinitely many at $L_4$). What is universal is the \emph{shape}: at every
level, $\sigma$-action produces a fixed/paired decomposition that is
structurally recognisable as the same kind of object.

\subsection{Why this is non-trivial --- and where the non-triviality currently rests on hypothesis}\label{ssec:nontrivial}

\begin{remark}[On the non-triviality of the universal pattern, and its load-bearing conditions]
A skeptic might respond: \emph{any} involution on \emph{any} set produces
fixed and paired subsets. This is trivially true at the set-theoretic
level. The universal-$\sigma$-pattern claim of
Conditional~\ref{thm:universal} is non-trivial only if three load-bearing
conditions hold; each is named here:
\begin{enumerate}\itemsep0pt
  \item \textbf{H-SigmaRestrict:} the involution at every level is the
    \emph{same} involution in different coordinates --- the cascade
    $\sigma$, which restricts to $\sigma_K$ on $\mathbb{Q}(\sqrt{5})$
    (Conditional~\ref{lem:sigma_restricts}).
  \item \textbf{H-UniqueSigma:} at each level, the involution is the
    unique one naturally associated with $\mathcal{C}_\varphi$ on that
    level's objects, not an arbitrary choice. (The skeptic's worry is
    that some \emph{other} involution might also produce a fixed/paired
    decomposition at some level by coincidence; H-UniqueSigma asserts
    that the cascade $\sigma$ is structurally privileged.)
  \item \textbf{H-NaturalIrreducibles:} the objects at each level
    (vertices, orbits, eigenspaces, $K$-poles, primes) are the natural
    closure-irreducibles at that level, not artificial constructs
    selected post-hoc to make the pattern hold.
\end{enumerate}
Under all three hypotheses, the non-triviality is: \emph{the same
closure-derived involution $\sigma$ acts at every level of the typology,
and at every level the fixed/paired decomposition is the natural
cascade-respecting classification.} This is a structural fact, not a
tautology --- but \emph{it depends on the three hypotheses above}.

\smallskip\noindent\emph{Honest scope of this paper.} We do not prove
H-UniqueSigma or H-NaturalIrreducibles. We name them so a reader can see
where the non-triviality sits. A reader who rejects either hypothesis is
free to read Conditional~\ref{thm:universal} as a coincidental
combinatorial pattern; a reader who accepts them gets a structural
claim. The pattern itself --- that every level admits a fixed/paired
decomposition --- is unconditional. What is conditional is the claim
that this pattern is non-tautological.

The non-triviality is therefore: the same closure-derived involution
$\sigma$ acts at every level of the typology, and at every level the
fixed/paired decomposition is the natural cascade-respecting
classification. This is not tautological; it is a structural fact about
how the cascade $\sigma$-action propagates through all levels of the
closure-irreducibility hierarchy.
\end{remark}

\section{What this bridge does and does not provide}\label{sec:does_does_not}

\subsection{What the $L_4$ placement provides}

\begin{itemize}\itemsep0pt
  \item \textbf{Structural placement of classical primes.} Every classical prime has a well-defined cascade-$\sigma$-class. No prime is foreign to the typology.
  \item \textbf{Population of the typology.} The five-level structure now has all five levels populated with concrete object sets; $L_4$ is developed via the Dedekind splitting law. We do not claim that $L_0$--$L_4$ exhaust all conceivable closure-irreducibles; we claim that these five levels are demonstrably populated, and that classical primes specifically belong to $L_4$.
  \item \textbf{Demonstration of universal $\sigma$-pattern.} The same closure-derived involution acts at every level, producing the same shape of classification.
  \item \textbf{A foundational answer to ``what is a prime?''} A prime is a closure-irreducible at some level $L$; classical integer primes are the $L_4$ instance.
  \item \textbf{Bridge to the existence programme.} Via Paper~I's observation-manifold theorem, the closure framework's identification of $\sigma$-fixedness with bounded reference frames (\cite{SmartExistenceClosure}) extends to a structural reading of integer primes as $L_4$ closure-irreducibles.
\end{itemize}

\subsection{What it does not provide}

\begin{itemize}\itemsep0pt
  \item \textbf{Derivation of classical primes.} The integers $2, 3, 5, 7, \ldots$ are not derived from cascade; they exist as facts of arithmetic, and we classify them post-hoc within the cascade typology.
  \item \textbf{Derivation of the Euler product.} The classical $\zeta(s) = \prod_p (1 - p^{-s})^{-1}$ is not obtained from cascade machinery; the Hecke obstruction (rh-formal $M3$) blocks identification of $\hat{\zeta}$ with $\zeta$.
  \item \textbf{Proof of classical RH.} The arithmetic question of whether classical $\zeta$ zeros lie on the critical line is not addressed; it remains as classical RH.
  \item \textbf{Derivation of the prime number theorem.} The asymptotic distribution $\pi(N) \sim N / \log N$ is not derived; only the asymptotic density within each $\sigma$-class (split / inert at $1/2$ each, ramified at $0$) is recorded.
\end{itemize}

\subsection{Open conditional propositions}

\begin{conditional}[H-EulerBridge]\label{cond:euler}
There exists a cascade-natural construction (operator, distribution, or
functional) producing a generating function $G(s)$ whose Euler-style
factorisation reduces, under a yet-to-be-specified completion process,
to the classical Euler product $\zeta(s) = \prod_p (1 - p^{-s})^{-1}$
of the $L_4$ primes. Under H-EulerBridge, the bridge from cascade $\hat{\zeta}$
to classical $\zeta$ is no longer obstructed at the multiplicity level.
\emph{Status: open. The Hecke obstruction (rh-formal $M3$) is the
specific named barrier.}
\end{conditional}

\begin{conditional}[H-PrimeDistribution]\label{cond:distribution}
Cascade dynamics on $V_{600}$ produces an asymptotic counting law for
$L_4$ closure-irreducibles (classical primes) that, under
H-EulerBridge or a related bridge condition, reduces to the classical
prime-counting function $\pi(N)$. \emph{Status: open.}
\end{conditional}

These two conditional propositions are stated for completeness; they are
not load-bearing for the present paper, whose main claims
(Theorems~\ref{thm:dedekind}, \ref{thm:L4_placement},
\ref{thm:universal}) are all unconditional given Paper~I.

\section{The Hecke obstruction: acknowledged, not resolved}\label{sec:hecke}

A careful reader will ask: \emph{if cascade $\sigma$ classifies classical
primes via the splitting law, does this mean the cascade analytic object
$\hat{\zeta}(z)$ corresponds to classical $\zeta(s)$ in some way?}

The answer is \emph{no, currently}, and we record why.

\subsection{The cascade analytic object $\hat{\zeta}(z)$}

Per rh-formal (\cite{rhFormalSmart} $M1$):
\[
\hat{\zeta}(z) = \prod_{K \in \{0, 20, 52, 72\}} \left( 1 - L_K z^{10} + z^{20} \right)^{-m_K}
\]
with multiplicities $(m_0, m_{20}, m_{52}, m_{72}) = (1, 1, 5, 5)$ and
$L_K$ values constructed from $\varphi$ via the pentagonal-clock
structure.

\subsection{Why $\hat{\zeta} \neq \zeta$}

Three independent mismatches:

\begin{enumerate}\itemsep0pt
  \item \textbf{Factor count.} $\hat{\zeta}$ has four Euler factors (one per $K$-value); classical $\zeta$ has infinitely many (one per prime).
  \item \textbf{Multiplicity pattern.} $\hat{\zeta}$ has multiplicities $(1, 1, 5, 5)$; classical $\zeta$ has multiplicities all $1$.
  \item \textbf{Coefficients.} $L_K$ are golden-ratio expressions; classical Hecke eigenvalues for $\zeta$ are all $1$ (or for general L-functions, follow Ramanujan-style bounds $|a_p| \le 2\sqrt{p}$ which $L_K$ do not satisfy in the cascade frame).
\end{enumerate}

\emph{These are not framing problems.} The multiplicities $(1, 1, 5, 5)$
are counted geometric objects on $V_{600}$ under pentagonal-clock
$\sigma$-action. They are what the cascade produces. No reinterpretation
makes them match classical Hecke arithmetic.

\subsection{What this means for the bridge}

The $L_4$ $\sigma$-classification (this paper) is independent of the
Hecke obstruction. The splitting law classifies primes; it does not
require identifying $\hat{\zeta}$ with $\zeta$. The placement of
classical primes in the cascade typology is therefore unconditional and
does not depend on resolving the Hecke obstruction.

The Hecke obstruction would only need to be resolved if one wanted to:

\begin{itemize}\itemsep0pt
  \item Derive classical $\zeta$ from cascade closure (Conditional Proposition~\ref{cond:euler}).
  \item Prove classical RH via cascade.
  \item Identify $\hat{\zeta}$ as ``the same object'' as $\zeta$.
\end{itemize}

This paper does \emph{none} of these. The $L_4$ classification is a
weaker, structural claim --- and is fully sufficient for its stated
purpose (placement of classical primes within the typology).

\section{Falsifiers and limitations}\label{sec:falsifiers}

\paragraph{Structural falsifiers (load-bearing).}
\begin{enumerate}\itemsep0pt
  \item \textbf{F-L4-1.} A constructed counterexample showing the cascade $\sigma$ does not restrict to $\sigma_K$ on $\mathbb{Q}(\sqrt{5})$ falsifies Lemma~\ref{lem:sigma_restricts} and the central identification of this paper.
  \item \textbf{F-L4-2.} A misstatement of the classical Dedekind splitting law in $\mathbb{Z}[\varphi]$ (Theorem~\ref{thm:dedekind}) would falsify the foundation; but this is textbook classical math (Dedekind 1871) and is not in dispute.
  \item \textbf{F-Universal-1.} A demonstration that one of the five levels $L_0$--$L_4$ does \emph{not} admit a natural $\sigma$-classification of the fixed/paired/special shape (e.g., shows a level where the $\sigma$-action is trivial or where the fixed/paired structure does not exist) falsifies the universal $\sigma$-pattern of Theorem~\ref{thm:universal}.
\end{enumerate}

\paragraph{Out-of-scope (does not falsify this paper).}
\begin{itemize}\itemsep0pt
  \item A proof of classical RH would not affect this paper; it would simply settle the $L_4$ arithmetic question that we explicitly do not address.
  \item A disproof of classical RH would also not affect this paper; the $\sigma$-classification of primes is independent of where $\zeta$ zeros sit.
  \item The Hecke obstruction (rh-formal $M3$) is acknowledged here; resolving or not resolving it does not affect the $L_4$ classification.
\end{itemize}

\paragraph{Acknowledged limitations.}
\begin{itemize}\itemsep0pt
  \item This paper is a \emph{structural placement}, not a derivation. Classical primes exist as facts of arithmetic; cascade does not produce them.
  \item The universal $\sigma$-pattern (Theorem~\ref{thm:universal}) records that the same shape appears at every level, but does not assert that the cascade $\sigma$ uniquely determines the $\sigma$-classification at each level. Other involutions may produce the same classification at some levels by coincidence; we do not claim cascade-$\sigma$-uniqueness.
  \item The Dedekind splitting law applies to $\mathbb{Q}(\sqrt{5})$ specifically. Other quadratic extensions have analogous splitting laws but a different fixed quadratic discriminant. The cascade's natural field is $\mathbb{Q}(\sqrt{5})$ by the icosian-quaternion construction (Section~\ref{sec:Q_sqrt5}); other choices are not natural to cascade and not considered here.
\end{itemize}

\section{Conclusion}\label{sec:conclusion}

We have shown that classical integer primes ($L_4$) fit into the
closure-irreducibility typology of Paper~I via the classical Dedekind
splitting law in $\mathbb{Z}[\varphi]$. Every classical prime has a
precise cascade-$\sigma$-class --- split, inert, or ramified ---
determined unconditionally by its behaviour mod $5$. The cascade $\sigma$
restricted to $\mathbb{Q}(\sqrt{5})$ is the Galois conjugation
$\sqrt{5} \mapsto -\sqrt{5}$, and the splitting law is the manifestation
of this $\sigma$-action at the integer-prime level.

The headline structural fact: \emph{the $\sigma$-pattern is universal
across all five levels of the typology}. The same closure-derived
involution acts at $L_0$ vertices, $L_1$ $\sigma$-orbits, $L_2$ spectral
classes, $L_3$ pentagonal-clock $K$-poles, and $L_4$ classical primes.
Each level admits a fixed / paired / (sometimes) ramified decomposition.
The closure-irreducibility framework is therefore complete at the
classification level.

What the present paper does \emph{not} provide is also clear: we do not
derive classical primes from cascade, we do not bridge the cascade
analytic object $\hat{\zeta}$ to classical $\zeta$ at the generating-function
level (the Hecke obstruction remains), and we do not prove or address
classical RH.

The contribution is structural placement: \emph{classical primes are one
specific level of a broader family of closure-irreducibles, and they
exhibit the same $\sigma$-pattern as every other level}. Together with
Paper~I, the closure-foundations programme now provides a complete
typology of primes-as-closure-irreducibles and a foundational account of
why the critical line exists as the observation manifold of the entire
structure.

\section{Reproduction}\label{sec:reproduction}

The classical Dedekind splitting law of Section~\ref{sec:dedekind}
and the $L_4$ classification of Section~\ref{sec:L4} are sim-verified
by the script \texttt{repro/verify\_dedekind\_splitting.py} accompanying
this paper (numpy only, deterministic, $\sim$300 lines). The script:
\begin{itemize}
\item sieves all rational primes $p \leq 10000$ and classifies each
      by behaviour in $\mathbb{Z}[\varphi]$ using the Dedekind criterion
      $\left(\tfrac{5}{p}\right)$ (equivalently, by $p \bmod 5$);
\item for the first 12 split primes, brute-force searches integer
      pairs $(a,b)$ with $a^2 + ab - b^2 = \pm p$ to exhibit an
      explicit factorisation $\pi = a + b\varphi$ together with its
      Galois conjugate $\sigma_K(\pi) = (a+b) - b\varphi$;
\item verifies the asymptotic density: split primes ($p \bmod 5 \in
      \{1,4\}$) $\approx 0.496$ and inert primes ($p \bmod 5 \in
      \{2,3\}$) $\approx 0.504$ at $P = 10000$ (Chebotarev / Dirichlet
      expectation $1/2$);
\item verifies that $\{5\}$ is the unique ramified prime;
\item builds the $L_0$--$L_4$ universal-$\sigma$-pattern table
      (Section~\ref{sec:universal}) with $L_0$--$L_3$ entries cited
      from Paper~I \cite{CriticalLinePaperI} and $L_4$ computed here.
\end{itemize}

The script is deterministic and produces JSON outputs
(\texttt{splitting\_classification.json},
\texttt{explicit\_factorisations.json},
\texttt{universal\_pattern\_table.json}) and two figures
(density-vs-$P$ and the universal pattern bar chart). All classical
checks pass; the universal-pattern recurrence across $L_0$--$L_4$
is exhibited, with the load-bearing identifications (H-SigmaRestrict,
H-UniqueSigma, H-NaturalIrreducibles) named in
Section~\ref{sec:universal} and not assumed in the sim itself.

\paragraph{Acknowledgements.}
This paper builds on classical algebraic number theory (Dedekind,
Kronecker, Hecke, and the subsequent twentieth-century development of
algebraic-number-theory textbooks), on Paper~I of the present series, and
on the broader Existence / Life / Closure programme. The author thanks
the classical number-theory community for the splitting-law foundation,
which has been textbook material for over a century and is here
reinterpreted within the closure framework.

\bibliographystyle{plainnat}
\bibliography{references}

\end{document}
